\chapter{La pipeline di elaborazione}
\label{ch:pipeline}

Tutti i modelli condividono la stessa sequenza di elaborazione. Le differenze fra un approccio e
l'altro riguardano il contenuto del passo~3 e, per i primi due modelli, la presenza del passo~1.

\section{I quattro passi}

\subsection*{Passo 1 --- Baseline degli apparecchi sempre accesi}

Frigoriferi e congelatori non hanno cicli di accensione riconoscibili: assorbono in modo
continuo, con un ciclo di compressore che il campionamento a 15 minuti non risolve. Nei primi due
modelli (\code{cvxpy} e \code{cvxpy\_activation}) il loro contributo viene stimato a monte con una
procedura dedicata e \textbf{sottratto} dal segnale; il solver lavora quindi su un residuo che
contiene i soli apparecchi a ciclo.

A partire da \code{cvxpy\_full} questo passo \textbf{non esiste più}: gli always-on entrano
direttamente nel modello come variabili, e il solver ottimizza simultaneamente il loro livello di
assorbimento e le attivazioni degli altri apparecchi. Il vantaggio è che i due problemi non sono
indipendenti --- una baseline stimata male distorce tutto ciò che segue --- e risolverli insieme
elimina l'errore di propagazione.

\subsection*{Passo 2 --- Ricampionamento}

Il segnale a un minuto viene aggregato a slot di ampiezza $g$ (default 15 minuti) usando una
funzione di aggregazione configurabile (media, massimo, minimo o mediana; il default è la media).

La granularità è il principale parametro di costo computazionale: il numero di variabili binarie
cresce linearmente con $T$, e il tempo di risoluzione di un MILP cresce molto più rapidamente.
Passare da 15 a 30 minuti dimezza $T$ e rende il problema sensibilmente più facile, al prezzo di
non poter più distinguere eventi di durata inferiore alla mezz'ora.

\subsection*{Passo 3 --- Ottimizzazione giorno per giorno}

Per ciascun giorno si costruisce e si risolve un MILP \textbf{indipendente}, con un budget di
tempo di \code{time\_limit} secondi (default 60).

La scelta di separare i giorni è dettata dalla complessità: un modello su 15 giorni avrebbe 15
volte le variabili e sarebbe fuori portata. Ha però una conseguenza strutturale importante,
discussa nel capitolo~\ref{ch:modelli}: \textbf{qualunque grandezza che debba propagarsi da un
giorno all'altro non può essere una variabile del modello}. È il caso della quota settimanale di
accensioni, che deve essere passata come costante e aggiornata all'esterno del solver.

\subsection*{Passo 4 --- Ricostruzione alla risoluzione originale}

La soluzione, definita sugli slot da $g$ minuti, viene riportata sull'indice a un minuto per
propagazione in avanti dell'ultimo valore noto, così da poter essere confrontata con il segnale
originale e sommata sull'intero periodo per il calcolo dei consumi in kWh.

\section{Selezione del periodo di analisi}

I modelli non vengono applicati all'intero storico disponibile ma a un campione di
\textbf{15 giorni per abitazione}, organizzati in 3 blocchi da 5 giorni consecutivi.

\begin{center}
\begin{tabularx}{\textwidth}{@{}l X@{}}
\toprule
\textbf{Criterio} & \textbf{Motivazione} \\
\midrule
Blocchi identici per tutte le abitazioni
& I blocchi sono estratti dai giorni presenti in \emph{tutti} i segnali, così le differenze fra
  case riflettono la casa e non il periodo osservato. \\
\addlinespace
Estrazione casuale ma riproducibile
& L'estrazione usa un seme fisso: rieseguire l'analisi produce gli stessi blocchi. \\
\addlinespace
Esclusione delle festività
& Nessun blocco può cadere fra il 22 dicembre e il 6 gennaio: i consumi di Natale, Capodanno ed
  Epifania non sono rappresentativi dell'uso ordinario. \\
\addlinespace
Blocchi distanziati
& Si impone una distanza minima fra gli inizi dei blocchi, per evitare che si addensino nello
  stesso periodo e per campionare condizioni diverse. \\
\addlinespace
Blocchi risolti separatamente
& I risultati vengono concatenati solo alla fine. Concatenare i segnali \emph{prima} del
  ricampionamento farebbe generare migliaia di slot vuoti negli intervalli fra un blocco e
  l'altro. \\
\bottomrule
\end{tabularx}
\end{center}
